\begin{figure}[tb]
\centering
\begin{tikzpicture}[>=Latex,font=\small,
 block/.style={draw,rounded corners=2pt,minimum height=1.2cm,align=center}]
 \node[block,fill=blue!5,minimum width=3.2cm] (m) at (0,0)
   {Original fibre directions\\$\widetilde M$, dimension $n$};
 \node[block,fill=orange!7,minimum width=3.3cm] (a) at (3.65,0)
   {Parameter directions\\$A$, dimension $q-n$};
 \node[block,fill=purple!5,minimum width=3.0cm] (t) at (7.2,0)
   {Stabilizing circles\\$\T^{kq}$, dimension $kq$};
 \draw[thick] (-1.6,-.8)--(-1.6,-.95)--(5.3,-.95)--(5.3,-.8);
 \node[align=center] (z) at (1.85,-1.45)
   {Base $Z$, dimension $q$};
 \draw[->] (z.south)--node[right] {$F_k$, nonzero degree} (1.85,-2.65);
 \node at (1.85,-2.9) {$S^q$};
 \node[align=center] at (7.2,-1.6)
   {Scalar curvature is controlled\\on $Z\times\T^{kq}$};
\end{tikzpicture}
\caption{The local direction splitting after $k$ finite passes. The base
$Z$ is locally a product of original and parameter directions; no global
product decomposition is asserted. The inverse-metric trace uses only the
$n$ original directions. The sphere map uses the full base of dimension
$q$, which is fixed before the number $kq$ of circles is chosen.}
\label{fig:direction-splitting}
\end{figure}
